% Generated from locale/tr/content/first-order-logic/introduction/soundness-completeness.tex; do not hand-edit.


\olfileid[tr]{fol}{int}{scp}

\olsection{Sağlamlık ve Tamlık}

Birinci dereceden mantık için !!{derivation} sistemleri de tanıtacağız.
Mantıkçıların geliştirdiği birçok !!{derivation} sistemi vardır, fakat
hepsi !!{sentence}s arasında aynı !!{derivability} bağıntısını
tanımlar. Kesin olarak tanımlanmış belirli türde !!a{derivation}
varsa,~$\Gamma$'nın~$!A$'yı \emph{türettiğini},
$\Gamma\Proves !A$, söyleriz. !!^{derivation}s her zaman simgelerin
sonlu düzenlemeleridir: bir !!{sentence}s listesi ya da daha karmaşık
bir yapı olabilir. !!{derivation} sistemlerinin amacı, !!a{sentence}
ifadesinin bir~$\Gamma$ kümesi tarafından mantıksal olarak gerektirilip
gerektirilmediğini belirlemeye yarayan araç sağlamaktır. Bu amaca
hizmet edebilmeleri için $\Gamma\Entails !A$ ancak ve ancak
$\Gamma\Proves !A$ ise geçerli olmalıdır.

$\Gamma\Proves !A$ olduğu halde $\Gamma\Entails !A$ değilse
!!{derivation} sistemimiz fazla güçlüdür; gereğinden fazlasını
ispatlar. $\Gamma\Proves !A$ ise $\Gamma\Entails !A$ olmasına
\emph{sağlamlık} denir ve bu, iyi her !!{derivation} sistemi için
asgari bir gerekliliktir. Öte yandan $\Gamma\Entails !A$ olduğu halde
$\Gamma\Proves !A$ değilse !!{derivation} sistemimiz fazla zayıftır;
yeterince ispat vermez. $\Gamma\Entails !A$ ise $\Gamma\Proves !A$
olmasına \emph{tamlık} denir. Sağlamlığı ispatlamak genellikle görece
kolaydır (tümevarımlı olarak tanımlanan !!{derivation}s tek tek ele
alındığında bunların yapısı üzerinde tümevarımla). Tamlığı ispatlamak
daha zordur.

Sağlamlık ve tamlığın bir dizi önemli sonucu vardır. Bir
!!{sentence}s kümesi~$\Gamma$, bir çelişki (örneğin
$!A\land\lnot !A$) türetiyorsa \emph{tutarsız} olarak adlandırılır.
Tutarsız~$\Gamma$ kümelerinin hiçbir modeli olamaz; bunlar
sağlanamazdır. Tamlıktan tersi çıkar: tutarsız olmayan---ya da
diyeceğimiz gibi \emph{tutarlı}---her~$\Gamma$ kümesinin bir modeli
vardır. Aslında bu tamlığa denktir ve tamlığın gerçekten ispatlayacağımız
biçimidir. Bu derin ve belki şaşırtıcı bir sonuçtur:~$\Gamma$'dan
$!A\land\lnot !A$ ispatlayamamanız bile~$\Gamma$'nın betimlediği gibi
!!a{structure} bulunduğunu güvence altına alır. Böylece tamlık şu
soruya yanıt verir: Hangi !!{sentence}s kümelerinin modelleri vardır?
Yanıt: Tam olarak ve yalnızca tutarlı kümelerin.

Sağlamlık ve tamlık teoremlerinin iki önemli sonucu vardır: tıkızlık
teoremi ile Löwenheim--Skolem teoremi. Bunlar model kuramının önemli
sonuçlarıdır ve birçok ilginç sonucu ortaya koymak için kullanılabilir.
İkisine daha önce değindik: Birinci dereceden mantık, !!a{structure}
ifadesinin !!{domain} boyutunun sonlu ya da !!{nonenumerable} olduğunu
ifade edemez.

Tarihsel olarak bütün bunları---birinci dereceden mantığın sözdizim ve
semantiğinin nasıl tanımlanacağını, iyi !!{derivation} sistemlerinin
nasıl tanımlanacağını, bunların sağlam ve tam olduğunun nasıl
ispatlanacağını, birinci dereceden dillerde nelerin ifade edilip
edilemeyeceğinin açıklığa kavuşturulmasını---çözmek ve doğru yapmak uzun
zaman aldı. Artık nasıl yapılacağını biliyoruz; yine de bütün
ayrıntılardan geçmek kafa karıştırıcı ve yorucu olabilir. Fakat aynı
zamanda önemlidir; çünkü birinci dereceden mantığın biçimsel dili için
burada geliştirilen yöntemler mantık, bilgisayar bilimi ve dilbilimin
her yanında uygulanır. Bu yüzden ayrıntıları çalışmak uzun vadede
karşılığını verir.

